\begin{enumerate}
    \item Puisque les forces de frottement sont négligeables, le bilan des
          forces se réduit à~:
          \begin{itemize}
              \item la force de traction $\vec{T}$ de direction verticale, de
                    sens vers le haut,
                    \hfill\textbf{(0,25pt)}
              \item le poids de l'ascenseur $\vec{P}$ de direction verticale, de
                    sens vers le bas et d'intensité
                    $P=m\times g=850\times 9,80=\qty{8.33e3}{\N}$.
                    \hfill\textbf{(0,25pt)}
          \end{itemize}
    \item Puisqu'il est dit que l'ascenseur monte à vitesse constante,
          \hfill\textbf{(0,5pt)} \\
          nous pouvons appliquer le principe d'inertie.
          \hfill\textbf{(0,5pt)}
    \item D'après le principe d'inertie, on peut affirmer que
          $\sum\vec{F}_{\text{ext}}=\vec{T}+\vec{P}=\vec{0}$.
          \hfill\textbf{(0,5pt)} \\
          Ceci veut donc dire que
          $\vec{T}=-\vec{P}$ et que $T=P=\qty{8.33e3}{\N}$.
          \hfill\textbf{(0,5pt)}
    \item $W_{AB}(\vec{P})=\vec{P}\cdot\vv{AB}=W_{AB}(\vec{P})=
          P\times AB\times\cos(\ang{180})$
          \hfill\textbf{(0,25pt)}

          $W_{AB}(\vec{P})=-\num{8.33e3}\times 60,0=\qty{-5.00e5}{\J}$.
          \hfill\textbf{(0,25pt)}

          Ce travail est \textbf{négatif} donc \textbf{résistant}.
          \hfill\textbf{(0,5pt)}
    \item La force de traction étant opposée au poids, son travail est moteur et
          vaut l'inverse du travail du poids.
          
          $\mathcal{P}_{m}=\frac{W_{AB}(\vec{T})}{\Delta t}=
          -\frac{W_{AB}(\vec{P})}{\Delta t}$
          \hfill\textbf{(0,5pt)}

          $\mathcal{P}_{m}=-\frac{\bigl(\num{-5.00e5}\bigr)}{0,5\times 60}=
          \qty{1.7e4}{\W}$.
          \hfill\textbf{(0,5pt)}
\end{enumerate}
